\section{Results}
\label{sec:results}

\subsection{Generic mixed-source harm was not confirmed}

The developmental factorial yielded 44/100 successes for FF, 40/100 for OO,
62/100 for FO, and 17/100 for OF (Fig.~\ref{fig:development}). Same-source
cells therefore averaged 42.0\% and mixed-source cells averaged 39.5\%. The
preregistered coherence contrast was only $+.025$. Its paired-block 95\%
interval was $[-.030,+.085]$, and its task-cluster interval was
$[-.005,+.055]$. Both include zero. This experiment did not confirm a generic
penalty for combining sources.

The directional pattern was analyzed after the registered cells were complete.
The post-hoc fresh-arm marginal effect was $+.245$, with paired-block and
task-cluster intervals of $[.155,.330]$ and $[.170,.330]$. The post-hoc
old-gripper effect was $+.205$, with intervals of $[.140,.275]$ and
$[.105,.305]$. FO exceeded OO on every task and was at least as successful as
FF on every task, improving on eight and tying on two. These observations
selected FO20 for confirmation; they are not an independent replication of it.

\begin{figure}[t]
    \centering
    \includegraphics[width=\columnwidth]{fig2_developmental_factorial.pdf}
    \caption{\textbf{Developmental temporal-source factorial.} Cell labels
    give successes from $n=100$ paired task-state blocks per condition. The
    registered coherence contrast compares the mean of FF and OO with the mean
    of FO and OF; its paired-block and task-cluster 95\% bootstrap intervals
    both span zero. The highlighted FO direction and the marginal source
    effects were identified post hoc and used only to define the later
    confirmatory hypothesis.}
    \label{fig:development}
\end{figure}

\subsection{Offline accuracy and closed-loop utility disagree}

At $d=20$, the teacher-forced audit preferred the old source for every measured
component (Fig.~\ref{fig:offline_closedloop}). Normalized translation loss was
.59578 for fresh and .50667 for old; normalized rotation loss was 1.12962 for
fresh and 1.09877 for old. Gripper sign error was .30760 for fresh and .27400
for old. Lower is better for all three metrics. Closed-loop development instead
favored a fresh arm and old gripper.

The disagreement is component-specific: the gripper direction matches across
offline and closed-loop views, while the arm direction reverses. It does not
show that demonstrations are wrong or identify a causal mechanism. It shows
that teacher-forced component prediction accuracy and closed-loop temporal-source
utility are not equivalent objectives in this audited setting. The observed
offline balanced interaction is numerical zero to floating-point tolerance, as
required by Eq.~\eqref{eq:separable_identity}.

\begin{figure}[t]
    \centering
    \includegraphics[width=\columnwidth]{fig3_offline_closed_loop.pdf}
    \caption{\textbf{Temporal-source preference depends on the measurement
    objective.} Arrows indicate the favored source rather than placing
    incomparable raw quantities on one axis. Teacher-forced losses over 82
    demonstration episodes and 10,654 eligible targets favor old arm and old
    gripper predictions. Post-hoc closed-loop marginal effects over 100 paired
    task-state blocks favor a fresh arm and old gripper. Lower error is better
    offline; higher success is better closed loop.}
    \label{fig:offline_closedloop}
\end{figure}

\subsection{Asymmetric temporal reuse is confirmed}

FO20 achieved 80/126 successes (63.5\%) on the primary cohort
(Fig.~\ref{fig:confirmation}a). Newest achieved 53/126 (42.1\%), full old20
55/126 (43.7\%), age-exponential ensembling 62/126 (49.2\%), and CogACT 59/126
(46.8\%). Thus FO20 improved over both single-source controls and both shared
full-action temporal-weighting baselines.

All four registered contrasts were stable positive
(Fig.~\ref{fig:confirmation}b). FO20 minus newest was $+.214$, with
paired-block CI $[+.143,+.294]$ and task-cluster CI $[+.127,+.310]$. FO20 minus
full old20 was $+.198$, with intervals $[+.095,+.302]$ and
$[+.063,+.317]$. FO20 minus age-exponential was $+.143$, with intervals
$[+.048,+.238]$ and $[+.016,+.278]$. FO20 minus CogACT was $+.167$, with
intervals $[+.071,+.262]$ and $[+.040,+.294]$. Every leave-one-task-out
estimate was positive for every comparison.

The aggregate result does not imply uniform task-level improvement. Against
full old20, for example, FO20 was lower on one task, tied on one, and higher on
seven. The task-cluster intervals and leave-one-task-out checks explicitly
carry this heterogeneity into the inference. The conclusion is therefore an
aggregate, task-spanning effect for the frozen primary cohort.

\begin{figure*}[t]
    \centering
    \includegraphics[width=0.38\textwidth]{fig4a_success_rates.pdf}\hfill
    \includegraphics[width=0.56\textwidth]{fig4b_contrasts.pdf}
    \caption{\textbf{Preregistered confirmation on untouched primary states.}
    \textbf{a}, Success counts and rates for the five executors over $n=126$
    paired task-state blocks (nine tasks, 14 states each). \textbf{b}, Four
    registered FO20-minus-comparator success-rate differences. Circles are
    paired estimates; thick and thin intervals are 95\% paired-block and
    task-cluster percentile bootstrap intervals, respectively, each from
    20,000 draws. The dashed line marks no difference. All interval lower
    bounds and every leave-one-task-out estimate are positive.}
    \label{fig:confirmation}
\end{figure*}
